\section{Conclusion}
\label{sec:conclusion}

In this work, we explored the viability of classic value-based Reinforcement Learning algorithms for high-dimensional, pixel-based continuous control. By systematically evaluating Branching DQN and Rainbow DQN across the DeepMind Control Suite, we demonstrated that standard Q-learning can successfully navigate complex walking tasks when paired with the right architectural and environmental adaptations.

A primary takeaway from our experiments is the absolute necessity of careful reward shaping. Standard linear velocity rewards often fail because they inadvertently encourage agents to maximize short-term gains through erratic, kamikaze-style forward lunges that result in immediate collapse. By replacing these with a smooth, exponential kernel-based reward, we forced the agent to learn stable and sustainable gaits. Additionally, meticulously tuning the ratio between forward propulsion incentives and healthy survival bonuses proved critical to balancing speed with mechanical stability.

Our results also highlight that successfully learning from raw pixels requires deeply tailored network architectures. The chosen observation resolution and the specific Convolutional Neural Network used for feature extraction must be capable of capturing the fine biomechanical details of the environment. Furthermore, standard network capacities often bottleneck performance in visual regimes. We found that significantly expanding the hidden dimensions of the network (such as upgrading the dueling head to 512 neurons) was strictly essential. Without this expanded capacity, the model cannot properly map extracted latent visual features to accurate action-value estimates.

Finally, applying value-based methods to physical environments demands an effective strategy for action space discretization. As the number of controllable joints increases, agents face a paralyzing combinatorial explosion of possible actions. Our findings reinforce that branching architectures are a highly effective solution to this problem. By independently estimating values for each joint, branching entirely bypasses the exponential growth of the action space, allowing value-based methods to scale gracefully to complex, multi-joint morphologies.

Ultimately, while policy gradient methods remain the default for continuous control, our research confirms that heavily augmented value-based algorithms like Rainbow DQN offer a highly sample-efficient and powerful alternative for visual-based walking tasks.